\documentclass[../../../include/open-logic-section]{subfiles}

\begin{document}

\olfileid[id]{sth}{z}{infinity-again}
\olsection{Ketakhinggaan}

Kita sudah memiliki cukup banyak aksioma untuk menjamin adanya tak berhingga
banyak himpunan (jika ada himpunan sama sekali). Andaikan suatu himpunan ada;
maka $\emptyset$ ada (berdasarkan
\olref[sth][z][sep]{prop:emptyexists}). Untuk setiap himpunan~$x$, himpunan
$x\cup\{x\}$ ada berdasarkan
\olref[sth][z][pairs]{prop:pairsconsequences}. Jadi, dengan menerapkan operasi
ini beberapa kali, kita memperoleh himpunan-himpunan berikut:
\begin{enumerate}
	\item[0.] $\emptyset$ %& $0$ & stage $0$\\
	\item[1.] $ \{\emptyset\}$ %& $1$ & stage $1$\\
	\item[2.] $\{\emptyset, \{\emptyset\}\}$ %& $2$ & stage $2$\\
	\item[3.] $\{\emptyset, \{\emptyset\}, \{\emptyset, \{\emptyset\}\}\}$ %&$3$ & stage $3$\\
	\item[4.] $\{\emptyset, \{\emptyset\}, \{\emptyset, \{\emptyset\}\}, \{\emptyset, \{\emptyset\}, \{\emptyset, \{\emptyset\}\}\}\}$% & $4$ & stage $4$ 
\end{enumerate}
dan kita dapat memeriksa bahwa setiap himpunan ini berbeda.

Kita memulai penomoran dari~$0$ karena beberapa alasan. Salah satunya adalah
sebagai berikut. Tidak sulit untuk memeriksa bahwa himpunan yang kita beri
label ``$n$'' mempunyai tepat $n$ anggota dan (secara intuitif) dibentuk pada
tahap ke-$n$.

Namun, konstruksi ini hanya memberikan \emph{tak berhingga banyak} himpunan;
konstruksi ini tidak menjamin adanya suatu \emph{himpunan tak hingga}, yakni
himpunan dengan tak berhingga banyak anggota. Hal ini sungguh penting: jika
kita tidak dapat menemukan suatu himpunan yang tak hingga Dedekind, kita tidak
dapat mengonstruksi aljabar Dedekind. Padahal kita menginginkan aljabar
Dedekind agar dapat memperlakukannya sebagai himpunan bilangan asli.
(Bandingkan \olref[sfr][infinite][dedekindsproof]{sec}.)

Yang penting, aksioma-aksioma yang telah kita tetapkan sejauh ini
\emph{tidak} menjamin keberadaan himpunan tak hingga apa pun. Jadi, kita harus
menetapkan suatu aksioma baru:

\begin{axiom}[Ketakhinggaan]
Terdapat himpunan~$I$ sedemikian sehingga $\emptyset\in I$ dan
$x\cup\{x\}\in I$ setiap kali $x\in I$.
\begin{align*}
	\exists I( & ((\exists o \in I)\forall x\ x \notin o) \land {}\\
	& (\forall x \in I)(\exists s \in I)\forall z(z \in s \liff (z \in x \lor z = x)))
\end{align*}
\end{axiom}

Mudah dilihat bahwa himpunan~$I$ yang diberikan oleh Aksioma Ketakhinggaan
bersifat tak hingga Dedekind. Anggota istimewanya ialah $\emptyset$, dan
injeksi pada~$I$ diberikan oleh $s(x)=x\cup\{x\}$. Teorema
\olref[sfr][infinite][dedekind]{thm:DedekindInfiniteAlgebra} menunjukkan cara
mengekstrak suatu aljabar Dedekind dari himpunan yang tak hingga Dedekind;
kita akan memperlakukan aljabar ini sebagai himpunan bilangan asli kita.
Secara lebih tepat:
\begin{defn}\ollabel{defnomega}
Misalkan $I$ sebarang himpunan yang diberikan oleh Aksioma Ketakhinggaan.
Misalkan $s$ fungsi yang didefinisikan oleh $s(x)=x\cup\{x\}$. Tetapkan
$\omega=\closureofunder{s}{\emptyset}$. Kita menyebut anggota-anggota
$\omega$ sebagai \emph{bilangan asli}, dan mengatakan bahwa $n$ merupakan
hasil dari menerapkan $s$ kepada $\emptyset$ sebanyak $n$ kali.
\end{defn}

Sekarang Anda dapat kembali memeriksa bahwa himpunan yang beberapa paragraf
sebelumnya diberi label ``$n$'' akan diperlakukan \emph{sebagai} bilangan~$n$.

Kita akan membahas arti penting ketetapan ini dalam
\olref[sth][z][nat]{sec}. Untuk saat ini, ketetapan tersebut memungkinkan
kita membuktikan hasil yang intuitif:

\begin{prop}\ollabel{naturalnumbersarentinfinite}
Tidak ada bilangan asli yang bersifat tak hingga Dedekind.
\end{prop}

\begin{proof}
Bukti dilakukan dengan induksi, yakni
\olref[sfr][infinite][induction]{thm:dedinfiniteinduction}. Jelas bahwa
$0=\emptyset$ tidak bersifat tak hingga Dedekind. Untuk langkah induksi, kita
akan membuktikan kontraposisinya: jika (secara absurd) $s(n)$ bersifat tak
hingga Dedekind, maka $n$ bersifat tak hingga Dedekind.

Andaikan $s(n)$ bersifat tak hingga Dedekind; dengan kata lain, terdapat suatu
!!{injection}~$f$ dengan $\ran{f}\subsetneq\dom{f}=s(n)=n\cup\{n\}$. Ada dua
kasus yang perlu dipertimbangkan.

\emph{Kasus 1: $n\notin\ran{f}$.} Maka $\ran{f}\subseteq n$ dan
$f(n)\in n$. Misalkan $g=\funrestrictionto{f}{n}$; sekarang
$\ran{g}=\ran{f}\setminus\{f(n)\}\subsetneq n=\dom{g}$. Jadi, $n$ bersifat
tak hingga Dedekind.

\emph{Kasus 2: $n\in\ran{f}$.} Tetapkan suatu
$m\in\dom{f}\setminus\ran{f}$, dan definisikan fungsi~$h$ dengan domain
$s(n)=n\cup\{n\}$:
\[
h(x) = 
	\begin{cases}
		f(x) & \text{jika }f(x) \neq n\\
		m & \text{jika }f(x)=n.
	\end{cases}
\]
Jadi, $h$ dan $f$ bernilai sama di semua tempat, kecuali bahwa
$h(f^{-1}(n))=m\neq n=f(f^{-1}(n))$. Karena $f$
merupakan !!a{injection}, berlaku $n\notin\ran{h}$; dan
$\ran{h}\subsetneq\dom{h}=s(n)$. Sekarang, dengan menggunakan argumen pada
Kasus~1, $n$ bersifat tak hingga Dedekind.
\end{proof}

Namun, masih tersisa pertanyaan tentang bagaimana kita dapat
\emph{membenarkan} Aksioma Ketakhinggaan. Jawaban singkatnya adalah bahwa kita
perlu menambahkan satu prinsip lagi pada kisah kita. Prinsip itu adalah
sebagai berikut:
\begin{enumerate}
	\item[] \stagesinf. Terdapat suatu tahap tak hingga. Artinya, terdapat
	tahap yang (a) bukan tahap pertama, (b) mempunyai beberapa tahap yang
	mendahuluinya, tetapi (c) tidak mempunyai pendahulu langsung.
\end{enumerate}
Aksioma Ketakhinggaan langsung mengikuti prinsip ini. Kita tahu bahwa
bilangan asli~$n$ dibentuk pada tahap~$n$. Jadi, himpunan $\omega$ dibentuk
pada tahap tak hingga pertama. $\omega$ sendiri menjadi saksi bagi Aksioma
Ketakhinggaan.

Namun, hal ini hanya mengembalikan kita pada pertanyaan tentang bagaimana
kita dapat membenarkan \stagesinf. Seperti halnya \stagessucc, prinsip ini
bukan bagian eksplisit dari kisah kita mengenai hierarki kumulatif-iteratif.
Lebih dari itu, tidak ada apa pun dalam gagasan hierarki iteratif itu
sendiri---tempat himpunan dibentuk tahap demi tahap---yang memaksa kita
berpikir bahwa proses tersebut melibatkan suatu tahap \emph{tak hingga}.
Sangat masuk akal untuk berpikir bahwa tahap-tahap itu ditata seperti
bilangan asli.

Namun, hal ini menimbulkan masalah yang jelas. Dalam
\olref[sfr][infinite][dedekindsproof]{sec}, kita mempertimbangkan ``bukti''
Dedekind bahwa terdapat suatu himpunan (pikiran-pikiran) yang tak hingga
Dedekind. Mungkin bukti itu tidak terasa sangat memuaskan. Akan tetapi, jika
\stagesinf{} tidak ``dipaksakan kepada kita'' oleh konsepsi iteratif tentang
himpunan (atau oleh ``hukum-hukum pemikiran''), kita masih tidak mempunyai
pembenaran intrinsik bagi klaim bahwa terdapat suatu himpunan yang tak hingga
Dedekind.

Tentu saja, masih banyak yang dapat dikatakan. Namun, semoga sekarang Anda
telah sampai pada titik untuk mulai memikirkan apa yang mungkin
\emph{diperlukan} guna membenarkan suatu aksioma (atau prinsip). Dalam
pembahasan selanjutnya, kita menerima \stagesinf{} begitu saja.

\end{document}
